% !TEX root = ../../main.tex
% Adapted from the author's linear-algebra book on October 3, 2026.
% Editorial clarification: assigning each input an output defines a function;
% injectivity additionally requires that distinct inputs have distinct outputs.
\section{Injective, surjective, and bijective maps}
\label{sec:appendix:map-properties}

\subsection{Functions, domains, and codomains}

Let \(f:A\to B\) be a function, also called a map. The set \(A\) is its
\emph{domain}, and \(B\) is its \emph{codomain}. Every element of \(A\)
has exactly one output in \(B\). This requirement alone does not say
whether every element of \(B\) is reached or whether two different
inputs can have the same output.

The \emph{image} of \(f\) consists of the outputs that are actually reached,
\[
  f(A)=\{f(a)\mid a\in A\}\subseteq B.
\]
If \(b=f(a)\), we call \(a\) a \emph{preimage} of \(b\). An element of
the codomain may have no preimage, one preimage, or several preimages.
The domain and codomain are part of the function, so they must be specified
when deciding whether a map has any of the following properties.
\index{function}\index{image of a function}\index{preimage}

\subsection{Surjectivity}

The map \(f\) is \emph{surjective}, or \emph{onto}, if every element
of the codomain is the output of at least one input. Equivalently,
\[
  \forall b\in B\quad \exists a\in A\quad f(a)=b.
\]
In terms of the image, this says \(f(A)=B\).
Different inputs may have the same output. In the middle panel of
\autoref{fig:recall:map-properties}, every codomain element receives
an arrow, and two inputs share an output.
\index{surjective map}

\subsection{Injectivity}

The map \(f\) is \emph{injective}, or \emph{one-to-one}, if distinct
inputs have distinct outputs. Equivalently, for any \(a_1,a_2\in A\),
\[
  f(a_1)=f(a_2)\quad\Longrightarrow\quad a_1=a_2.
\]
\marginnote{As a mnemonic, think of the \emph{in} in \emph{injective}
  as a reminder to check the \emph{inputs}. Different inputs must
  remain distinguishable through their outputs.}
Each codomain element has at most one input mapping to it. Some
codomain elements may have no such input, as the left panel of
\autoref{fig:recall:map-properties} shows.
\index{injective map}

\subsection{Bijectivity}

The map \(f\) is \emph{bijective} if it is both injective and surjective.
Every codomain element then has exactly one input mapping to it.
This gives a one-to-one correspondence between the entire domain
and codomain, as in the right panel of
\autoref{fig:recall:map-properties}. Bijectivity imposes both
conditions, so it is stronger than either condition alone.
For a bijection, the \emph{inverse function} \(f^{-1}:B\to A\) sends
each \(b\in B\) to its unique preimage. Thus
\[
  f^{-1}(f(a))=a\quad(a\in A),\qquad
  f(f^{-1}(b))=b\quad(b\in B).
\]
\index{bijective map}\index{inverse function}

\begin{figure*}[htbp]
  \centering
  % Adapted from the author's linear-algebra book as a full-width comparison.
\begin{tikzpicture}[
  x=1mm,y=1mm,font=\footnotesize,line width=0.45pt,
  every node/.style={inner sep=2pt},
  map arrow/.style={line width=0.35pt,-{Latex[length=1.6mm,width=1.1mm]},
    shorten <=2.5pt,shorten >=2.5pt}
]
  \foreach \origin/\paneltitle in {
    0/{Injective, not surjective},
    52/{Surjective, not injective},
    104/{Bijective}
  } {
    \begin{scope}[xshift=\origin mm]
      \node at (21.5,28) {\paneltitle};
      \draw[fill=black!5] (9,0) ellipse[x radius=7mm,y radius=17mm];
      \draw[fill=black!5] (34,0) ellipse[x radius=7mm,y radius=17mm];
      \node at (9,22) {\(A\)};
      \node at (34,22) {\(B\)};
    \end{scope}
  }
  \begin{scope}
    \foreach \y in {9,0,-9} {
      \fill (9,\y) circle[radius=1.6pt];
    }
    \foreach \y in {9,3,-3,-9} {
      \fill (34,\y) circle[radius=1.6pt];
    }
    \draw[map arrow] (9,9) -- (34,9);
    \draw[map arrow] (9,0) -- (34,3);
    \draw[map arrow] (9,-9) -- (34,-3);
  \end{scope}
  \begin{scope}[xshift=52mm]
    \foreach \y in {9,3,-3,-9} {
      \fill (9,\y) circle[radius=1.6pt];
    }
    \foreach \y in {9,0,-9} {
      \fill (34,\y) circle[radius=1.6pt];
    }
    \draw[map arrow] (9,9) -- (34,9);
    \draw[map arrow] (9,3) -- (34,0);
    \draw[map arrow] (9,-3) -- (34,0);
    \draw[map arrow] (9,-9) -- (34,-9);
  \end{scope}
  \begin{scope}[xshift=104mm]
    \foreach \y in {9,0,-9} {
      \fill (9,\y) circle[radius=1.6pt];
      \fill (34,\y) circle[radius=1.6pt];
      \draw[map arrow] (9,\y) -- (34,\y);
    }
  \end{scope}
\end{tikzpicture}

  \caption[Injective, surjective, and bijective maps.]{Three maps between
    finite sets. Each dot is an element. The left map is injective and leaves
    an output unused. The middle map is surjective, but two inputs share an
    output. The right map is bijective and reaches every output exactly
    once. The ellipses enclose the elements of each set.}
  \label{fig:recall:map-properties}
\end{figure*}

\subsection{Examples in Euclidean spaces}

The same definitions apply to maps between spaces such as \(\R^2\)
and \(\R^3\). The domain and codomain are part of each example.

\paragraph{An injective inclusion of a plane.}
Consider
\[
  Z:\R^2\to\R^3,\qquad Z(x,y)=(x,y,0).
\]
Different pairs retain their distinct first two coordinates, so \(Z\)
is injective. Its image is the \(xy\)-plane. It is not surjective onto
\(\R^3\), since an output with a nonzero third coordinate is never reached.
\autoref{fig:recall:plane-inclusion} shows a pair and its image in that
plane.
If we instead took the codomain to be the \(xy\)-plane itself, the same
assignment would be bijective. This is why the codomain matters.

\begin{figure}[htbp]
  \centering
  % Schematic inclusion of R^2 as the xy-plane in R^3.
\begin{tikzpicture}[
  x=1mm,y=1mm,font=\footnotesize,line width=0.45pt,
  every node/.style={inner sep=2pt},
  axis/.style={-{Latex[length=1.6mm,width=1.1mm]}},
  map arrow/.style={line width=0.35pt,-{Latex[length=1.6mm,width=1.1mm]},
    shorten <=2.5pt,shorten >=2.5pt}
]
  \node at (10,29) {\(\R^2\)};
  \node at (72,29) {\(\R^3\)};
  \draw[axis] (0,0) -- (24,0) node[right=4pt] {\(x\)};
  \draw[axis] (0,0) -- (0,22) node[above=4pt] {\(y\)};
  \draw[gray] (0,0) -- (12,8);
  \fill (12,8) circle[radius=1.6pt];
  \begin{scope}[shift={(52,0)},x={(8mm,-3mm)},y={(8mm,3mm)},z={(0mm,9mm)}]
    \filldraw[fill=black!10] (0,0,0) -- (2,0,0) -- (2,2,0)
      -- (0,2,0) -- cycle;
    \draw[axis] (0,0,0) -- (2.7,0,0) node[below=4pt] {\(x\)};
    \draw[axis] (0,0,0) -- (0,2.7,0) node[above right=4pt] {\(y\)};
    \draw[axis] (0,0,0) -- (0,0,2.4) node[above=4pt] {\(z\)};
    \draw[gray] (0,0,0) -- (1.5,1,0);
    \fill (1.5,1,0) circle[radius=1.6pt];
  \end{scope}
  \draw[map arrow] (12,8) .. controls (30,19) and (49,13) .. (72,-1.5);
\end{tikzpicture}

  \caption[Including a plane in three-dimensional space.]{Schematic
    inclusion \(Z:\R^2\to\R^3\). The shaded region
    depicts part of the \(xy\)-plane, which continues beyond the drawing.
    The marked input maps to a point in this plane. Distinct input
    coordinates remain distinct after inclusion.}
  \label{fig:recall:plane-inclusion}
\end{figure}

\paragraph{A surjective coordinate projection.}
Consider
\[
  U:\R^3\to\R^2,\qquad U(x,y,z)=(y,z).
\]
Every pair \((y,z)\) is attained, for instance from \((0,y,z)\), so
\(U\) is surjective. It is not injective, since changing \(x\) leaves
the output unchanged. \autoref{fig:recall:coordinate-projection}
shows inputs differing in the \(x\)-direction sharing an output.

\begin{figure}[htbp]
  \centering
  % Schematic coordinate projection R^3 -> R^2, (x,y,z) -> (y,z).
\begin{tikzpicture}[
  analysis diagram,x=1mm,y=1mm,
  every node/.style={inner sep=2pt},
  axis/.style={-{Latex[length=1.6mm,width=1.1mm]}},
  map arrow/.style={-{Latex[length=1.6mm,width=1.1mm]},
    shorten <=1.6pt,shorten >=1.6pt}
]
  \node at (10,29) {\(\R^3\)};
  \node at (68,29) {\(\R^2\)};
  \begin{scope}[x={(8mm,-3mm)},y={(8mm,3mm)},z={(0mm,9mm)}]
    \draw[axis] (0,0,0) -- (3,0,0) node[below=4pt] {\(x\)};
    \draw[axis] (0,0,0) -- (0,2.6,0) node[below right=4pt] {\(y\)};
    \draw[axis] (0,0,0) -- (0,0,2.4) node[above=4pt] {\(z\)};
    \draw[gray,dashed] (0,0.8,1.8) -- (2.7,0.8,1.8);
    \coordinate (projection-input-one) at (1,0.8,1.8);
    \coordinate (projection-input-two) at (2,0.8,1.8);
    \fill (projection-input-one) circle[radius=1.6pt];
    \fill (projection-input-two) circle[radius=1.6pt];
  \end{scope}
  \draw[axis] (56,0) -- (80,0) node[right=4pt] {\(y\)};
  \draw[axis] (56,0) -- (56,22) node[above=4pt] {\(z\)};
  \coordinate (projection-output) at (64,18);
  \fill (projection-output) circle[radius=1.6pt];
  \draw[gray] (56,0) -- (projection-output);
  \draw[map arrow] (projection-input-one)
    .. controls (32,22) and (46,22) .. (projection-output);
  \draw[map arrow] (projection-input-two)
    .. controls (38,13) and (48,15) .. (projection-output);
\end{tikzpicture}

  \caption[Forgetting one coordinate.]{Schematic projection
    \(U:\R^3\to\R^2\). The marked inputs
    lie on a line parallel to the \(x\)-axis and have the same \(y,z\)
    coordinates. They map to the same pair. Every pair in the codomain
    is attained.}
  \label{fig:recall:coordinate-projection}
\end{figure}

\paragraph{A bijective coordinate swap.}
Consider
\[
  B:\R^2\to\R^2,\qquad B(x,y)=(y,x).
\]
Each output \((a,b)\) has exactly one preimage, namely \((b,a)\).
Thus \(B\) is bijective. Swapping coordinates again recovers the
original input. \autoref{fig:recall:coordinate-swap} shows the
reflection across the line \(x=y\).

\begin{figure}[htbp]
  \centering
  % Schematic coordinate swap R^2 -> R^2, (x,y) -> (y,x).
\begin{tikzpicture}[
  x=1mm,y=1mm,font=\footnotesize,line width=0.45pt,
  every node/.style={inner sep=2pt},
  axis/.style={-{Latex[length=1.6mm,width=1.1mm]}},
  map arrow/.style={line width=0.35pt,-{Latex[length=1.6mm,width=1.1mm]},
    shorten <=2.5pt,shorten >=2.5pt}
]
  \node at (10,29) {\(\R^2\)};
  \node at (68,29) {\(\R^2\)};
  \foreach \origin in {0,56} {
    \begin{scope}[xshift=\origin mm]
      \draw[axis] (0,0) -- (24,0) node[right=4pt] {\(x\)};
      \draw[axis] (0,0) -- (0,22) node[above=4pt] {\(y\)};
      \draw[gray,dashed] (0,0) -- (20,20);
    \end{scope}
  }
  \draw[gray] (0,0) -- (16,8);
  \draw[gray] (56,0) -- (64,16);
  \fill (16,8) circle[radius=1.6pt];
  \fill (64,16) circle[radius=1.6pt];
  \draw[map arrow] (16,8) .. controls (32,18) and (48,18) .. (64,16);
\end{tikzpicture}

  \caption[Swapping two coordinates.]{Schematic coordinate swap in
    \(\R^2\). The dashed diagonal
    in each coordinate system is \(x=y\). The marked point and its
    image have exchanged coordinates. Each output has exactly one input.}
  \label{fig:recall:coordinate-swap}
\end{figure}
